
% 来源：EPUB part039.xhtml
\section{投资简单但不容易}

很多年前俞斌问我什么叫“会投资”，我的回答就是这样：和“会下棋”一样。我肯定会下棋，俞斌也就是会下棋而已，但确实不是一回事。（2010-07-08）

% 来源：EPUB part040.xhtml
\subsection{道需悟，术可学}

简单指的是原则——就是不懂不做，不容易指的是理解、搞懂生意不容易。（2013-03-26）

投资的难其实就是因为其简单。买股票就是买公司，所以不懂不碰。要搞懂企业的道理很简单，真懂非常难。（2020-10-11）

知不易行也不易。实际上就是，知道什么是对的事情不容易，把事情做对了也不容易。认为知易行难的很可能还没完全知。（2023-11-17）

从做对的事情的角度讲，投资里做对的事情就是：买股票就是买公司，就是买公司的生意，就是买公司未来的净现金流折现，这点其实不用看书，懂了就懂了，不懂看啥书帮助都不会太大。

从把事情做对的角度看，投资就是要看懂公司的生意，有些书也许是可以帮助你看懂一些生意的。巴菲特的股东信是涵盖了两者的，如果反复看还是没帮助的话，别的书估计帮助也不大。我本人是没看过啥书的，但做过企业对理解企业有很大帮助。（2018-10-08）

\BookQuestion{网友：感觉您是一个能将复杂问题简单化的人，而且异常简洁。不知是天性，还是后天历练出的本事？}

呵呵，没觉得自己有什么特别。如果有的话实际上是我在这上面花的时间比较多而已。简化一个复杂的问题绝对不是一件容易的事。我的很多“简单”的结论是花了很长很长时间得出来的，千万别以为我是一眼就看出来了。不过，“平常心”可以帮助人去找到事物的本源。（2010-03-22）

\BookQuestion{网友：懂巴菲特的理念不太难，坚持巴菲特的投资理念难于上青天。}

坚持不了的人就是不懂！这里懂的意思是骨子里相信，不是可以照着念但其实不知道点什么。其实价值投资是所谓投资的唯一一条路，不然你投的是什么？（2014-06-25）

道需悟，术可学。

道指的就是做对的事情，或者说不做不对的事情，或者说发现错了马上改——不管多大的代价都是最小的代价。

术指的是技术，对每个人来说都是需要学同时也需要练习的，非一朝一夕之功。（2013-03-24）

先有道然后术才有用哈，不然就是瞎忙。（2015-01-16）

其实巴菲特并没有教我们什么。如果我们本来不懂的话，是没有人可以教会的。（2011-03-07）

我从来没有系统地看过巴菲特的东西。我非常认同买股票就是买公司这个逻辑。对我来说，这个逻辑其实不是来自巴菲特，而是来自我自己的内心，只是后来发现巴菲特也是如此，而且干得很不错，从而坚定了自己的信心而已。所以，我更认为我喜欢巴菲特，是因为他是我的“同道中人”。

巴菲特股东信里讲的东西无非就两个：一个是做对的事情，就是去找好公司，关注公司长远的未来；另一个是如何把事情做对，就是如何找到好公司，如何看到公司的未来。

第一个问题其实就一句话，非常典型的“大道至简”，但非常难，绝大多数人在这点上过不去，所以才会过于关注市场。

第二个问题，是每个人的能力圈问题，做投资，每个人都该找自己能搞懂的好企业。巴菲特懂的我经常不懂，我懂的巴菲特也可能不懂。我大致看过一些巴菲特的股东信，也很认真地听过他在几个大学的问答式演讲。他的语言精练、诙谐、逻辑性极强，对我来说听起来就好像听音乐一样，非常舒服。对于能看懂他的东西的人来说，这些已经足够好了；对于那些看不懂的人来说，其实就算详细讲解也是没什么用的。

总而言之，在投资里，能说的巴菲特都说过了，我没见过谁能比他说得更简单明了！（2016-08-30）

道的东西确实难教，必须要靠自己悟。心中无道的人你怎么说也是没用的。大多数人都是不容易开放的，很难接受跟自己以前认识不同的东西。我有时候感觉很奇怪的是，如果是个已经很成功的人，不太接受新东西我也能理解，因为人家已经有了足够成功的经历和理解了。现实中我看到的情况刚好相反，那些不开放的人常常就是那些不太成功的人。举个例子，炒了30年股票（也许10年、20年）都不赚钱的人，当我说巴菲特时，人家居然会说……布拉布拉的。（2019-04-11）

\BookQuestion{网友：改变思维方式还是有可能的，譬如我。}

呵呵，那真要恭喜你。其实我的本意也不是说绝对不可能变，但确实很少见。不过，有些人觉得自己变了可能是重新发现了自己而已。（2011-12-31）

\BookQuestion{网友：是不是心中没有本分和平常心就悟不到你的回答？只有心中有这两样东西才能领会？}

可能还需要时间吧？我花20年悟明白的东西也许你需要悟15年，如果你悟的话。（2019-09-16）

心中有道的人也许被点一下会有所悟，但那也是你自己的造化哈。（2024-09-10）

想起一个小故事。很多年前，我和华以刚下过一盘围棋，有一块我感觉非常不舒服，复盘的时候我问他为什么我怎么下都觉得不对，然后他告诉我：其实你这里怎么下都一样的，因为你前面下错了。（2013-04-28）

\BookQuestion{网友：在2016年看到您的一句话，“发现错了马上改，不管多大的代价都是最小的代价”，于是我坚决把一只明知是错但还在犹豫的股票卖掉了。当时亏了20\%，放到现在则是亏60\%。做到这一点确实不容易。}

不这么做会更不容易的！（2018-01-23）

\BookQuestion{网友：不这么做不但损失惨重，而且还会错失好的机会，时间越长这一正一负的差额越大，损失越大。}

对，就是机会成本的意思。（2018-01-27）

\BookQuestion{网友：怎么才能知道是做对了还是做错了？有些事只能事后知道。}

呵呵，非常好的问题。我对做对的事情的看法是：发现错了马上改，不管多大的代价都是最小的代价。其实这个世界上有很多人（包括我在内），经常会明知不对，但由于各种原因而不肯改，结果错误变得更大。呵呵，我肯定犯过这种错误。没人不犯错，但知道这个原则的人犯错概率低，改得快。（2010-03-14）

\BookQuestion{网友：问题是开始怎么知道是对的事情呢？}

你肯定知道有些事情是错的。做对的事情的意思是，如果错了就改，马上改的代价最小。很多人是明知是错还要做，但错的事情往往有眼前利益，有诱惑。比如：骗人的生意、抽烟、赌博，等等。（2010-05-25）

\BookQuestion{网友：比如炒股票、做短线就是错的事情（玩玩而已那可以理解），做真正的价值投资就是对的事情。当然偶尔犯错也是可以理解的，但不要犯原则性错误。是这样吗？}

嗯，你的解释很对。（2010-03-14）

\BookQuestion{网友：你在奥马哈能学到什么？}

虽然我觉得自己算是很理解价值投资了，但对巴菲特理解的那种厚度还是很震惊。这次巴菲特在回答一个做私募股权投资（Private Equity）的小伙子的问题时，只说了：“私募股权投资是买入为了卖出。如果我是你，我会试着赚一些钱然后买一些好公司来持有。”如果是我，说的意思可能差不多，但好像不能做到这么简单。（2010-05-13）

说一些和投资没那么直接相关的事：80岁的巴菲特和86岁的芒格在会上一讲（包括回答问题）就是近6个小时，好像身体好得很。第二天上午我见到他们时，依然觉得他们很精神，一点也没看出疲惫的样子。我突然觉得价值投资可能是个可以让人长寿的好办法。5号去听芒格讲，以前还没去过呢。

顺便说下，《穷查理宝典：查理·芒格智慧箴言录》这本书不错，李录的序言写得也很好，我翻了一下，觉得自己该好好看看。（2010-05-03）

在此再次推荐一下《穷查理宝典》，查理对“实际上非常难”的东西很有办法。

对刚刚开始了解价值投资的人来说，李录写的序非常值得认真读几遍。

希望不要觉得这本书贵，因为值。我是这个月初去洛杉矶参加查理的股东大会时买的。110美元一本，买了两本（没有查理签名的是60美元一本），我觉得这可能是我一生中最好的一笔投资（其实我两年前就买过两本原文的，一直没看，所以真的有点贵）。（2010-05-25）

% 来源：EPUB part041.xhtml
\subsection{投资没有充分条件}

投资里最难的就是没有具有“充分条件”的投资，只能从概率的角度考虑。越了解的公司，决策的把握越大，但从来没有100\%的赢面。所以仓位只能根据自己的理解来决定。（2013-03-01）

基本上，投资里是无法找到充分条件的，甚至必要条件都很难找到。

举个例子，为什么说“好的生意模式、好的企业文化、好的管理团队、好价钱”也不是可以买的充分条件？因为这样的公司也可能犯致命的错误，只是具备这种条件的公司犯致命错误的概率比没有这种条件的公司低很多，康复能力强很多，所以投资这种公司成功的概率高很多而已。

投资里最难的地方就是找不到可以投的充分条件，如果有的话，那电脑就可以比人脑厉害了。

“投资界”总是有人试图找出一种不用懂企业就可以在股市上赚钱的办法，他们永远都不会找到的。如果能找到的话，这个办法应该早就找到了，那现在投资界赚钱最多的一定是数学家。

在美国，我确实见过很好的大学里有些教授依然在试图找到用数学和电脑就可以赚钱的办法，他们的逻辑推理的前提里有些假设是不对的，所以后面的结论不管多“高级”都不可能是正确的结果。有意思的是，这些教授大多数都是我曾经的同行——计量经济出身的（我曾经是人大计量经济学的研究生哈）。（2013-03-29）

\BookQuestion{网友：芒格说过，与巴菲特工作了这么多年，发现他这个人的优点之一是他总是自觉地从数学的排列与组合的角度思考问题。}

我想他这里说的排列组合大概是指概率的意思。总的来讲芒格大概是认为巴菲特总是能理性地思考问题，这点对长期价值投资是非常重要的。

投资并不需要高深的数学，但基本的概率概念还是需要的。如果没学过概率，或者没学懂的话，简单跟随就行了。比如，没必要自己去试扔硬币正反面的实验。又比如，有博友说他去赌场下了10万次注还是赢的，有概率概念的人一看就知道这哥们的记忆有问题。（2010-04-18）

安全边际应该指的就是你自己的了解度，投资里没有数学上100\%的东西。（2015-01-21）

赌一般指的是48\%赢面的下注， 90\%以上赢面的下注我一般称之为投资。（2023-07-21）

\BookQuestion{网友：有七八成把握就可以动手吗？}

投资没有充分条件。你的七八成把握指的是什么很重要。如果你指的是你对公司10年后的理解，那把握度就是你的安全边际。（2015-01-22）

% 来源：EPUB part042.xhtml
\subsection{至少85\%的人不适合投资}

\BookQuestion{网友：您觉得一毕业就直接投身到投资领域好吗？还是应该先有自己经营企业的经验才能真正理解企业？}

老实讲，我不知道什么人适合做投资。但我知道统计上大概80\%\( - \)90\%进入股市的人都是赔钱的。如果算上利息的话，赔钱的比例还要高些。许多人很想做投资的原因可能是认为投资的钱比较好赚，或来得比较快。作为既经营过企业又有投资经验的人来讲，我个人认为经营企业还是要比投资容易些。虽然这两者其实没有什么本质差别，但经营企业总是会在自己熟悉的领域，犯错的机会小，而投资却总是需要面临很多新的东西和不确定性，而且投资人会非常容易变成投机者，从而去冒不该冒的风险，而投机者要转化为真正的投资者则可能要长得多的时间。

投资和投机其实是很不同的游戏，但看起来又非常像。就像在澳门，开赌场的就是投资者，而赌客就是投机者一样。赌场之所以总有源源不断的客源，是因为总有赌客能赢钱，而赢钱的总是比较大声些。作为娱乐，赌点小钱无可非议，但赌身家就不对了。可我真是能见到好多在股场上赌身家的人啊。

以我个人的观点，其实什么人都可以做投资，只要你明白自己买的是什么，价值在哪里。投机需要的技巧可能要高很多，这是我不太懂的领域，也不打算学了，有空还是多陪陪家人或打几场高尔夫吧。

作为刚毕业的学生，马上投入投资领域也没啥不可，但对企业的理解自然会弱些。但只要你热爱你干的事情，又知道自己的弱点，慢慢学习总是会明白的。

即使是号称很有企业经验的本人也是在经受很多挫折之后才觉得自己对投资的理解比较好了。我问过巴菲特在投资中不可以做的事情是什么，他告诉我说：不做空，不借钱，最重要的是不要做不懂的东西。这些年，我在投资里亏掉的美金数以亿计，每一笔都是在违背巴菲特教导的情况下亏的，而赚到的大钱也都是在自己真正懂的地方赚的。作为刚出道的学生，书上的东西可能知道得很多，但融到骨子里还需要吃很多亏后才行。所以，如果你马上投入投资行业，最重要的是要保守啊，别因为一个错误就再也爬不起来了。这里唯一我可以保证的是，你肯定会犯错误的。

无论如何，投资是个非常有趣的工作，如果你真准备好了，那就来吧。你真的准备好了吗？（2010-02-03）

\BookQuestion{网友：什么样的人是价值投资者？}

我从一开始就是个价值投资者。我买股票时总是假设，如果我有足够多的钱的话我是否会把整个公司买下来。（2010-03-27）

\BookQuestion{网友：投资是要对企业了解，但现实中许多企业家并不懂得投资股票，难道他们还不了解企业吗？}

一般来讲，好的企业家都很容易变成好的价值投资者，只要他们能悟到买股票就是买公司，悟到这一点非常不简单。（2011-02-11）

\BookQuestion{网友：如何判断一个人成功，是靠运气还是实力？}

如果他是凭运气赚到的钱，大多数早晚会凭实力亏掉。如果一个人几十年来运气都很好，千万别觉得那是幸存者偏差。巴菲特讲的那个扔硬币实验就是个很好的例子。（2024-06-08）

\BookQuestion{网友：我觉得百万以内的资金根本不用进股票市场，直接投到自己或者身边人熟悉的小生意上更安全可控。我在现实世界里的回报比股票高多了，虽然要操很多心，但也学到很多知识。自己经营过一点小生意也对看报表有很大的帮助，将来资金量大了以后再买股票可能眼光更准。}

非常同意你的观点。本来投上市公司的钱就应该是闲钱，如果自己能找到生意，那当然应该投自己的生意。其实投上市公司的内在道理是一样的，差别是不用自己费力去经营而已，很多时候甚至还能请到比自己更能干的人去经营。记得很久以前在学校碰到一个老师问，自己不懂投资但需要管理自己的小小财富怎么办。我说你就都交给巴菲特吧，全部买上巴菲特的股票，怎么着都会比自己管好的，还不用操心。（2013-05-29）

\BookQuestion{网友：个人专职投资能否作为赖以生存以及实现财富梦想的方式？}

话糙理不糙，我觉得你不行！当一个人需要问别人自己可不可以靠投资赚钱的时候，多半应该是不行的。投资最重要的是要有独立思考的能力，要能够理性地思考。（2013-10-29）

\BookQuestion{网友：这个月茅台掉30个点，一般人的心脏受不了。}

承受不了这点跌幅的人们本来就不应该碰股市。再说，股市上历年统计显示大部分人都是亏钱的。这个统计其实很怪异，因为长期来说，股市一直是涨的（比如用10年段看，过去几十年的每个10年段股市可能都是涨的），但为什么80\%以上的人们会是亏损的呢？（2022-10-30）

\BookQuestion{网友：虽然资金都在茅台、腾讯、苹果上，但这些归根结底还是学长你的能力圈（比如我始终想不透彻，未来的苹果为什么一定可以和当年的索尼不一样）。十几年过去，增加了一些商业的常识，稀里糊涂赚了些钱，但是没有构建出自己可持续的能力圈，有时也挺焦虑的。}

如果你确认你一直在想但确实又想不太明白的话，最好的办法也许就是买标普500，然后该干嘛干嘛去？对很多人而言，花太多时间在“投资”上不值得，这和你的学历没关系。（2023-03-09）

\BookQuestion{网友：请教您对指数基金的看法？因为弄懂一个公司太难了，自己也要上班。巴菲特也说过指数基金最终会跑赢绝大多数股票投资者。}

好像巴菲特早就讲过，对大多数投资人而言，在指数便宜的时候投资指数是最好的办法。我同意。（2010-02-07）

巴菲特指的是美国的指数。我想， 85\%以上的投资者买指数的结果会更好，问题是这85\%指的都是谁呢？很少有人会认为自己属于85\%的，你能这么想说明你有智慧哈。（2013-04-03）

据说85\%\( - \)90\%的人在股市上都是亏钱的，另外可能还有5\%的人打平？这些人为什么不买标普500呢？长期来讲至少也是接近10\%的复合年回报吧（大致，自己没真的研究过）。问题是这85\%\( - \)90\%的人也都觉得自己能打败标普500？他们可能既不了解标普500，也不了解他们自己？这其实是个很有意思的问题，不过不是关于投资的。

对在股市上老是亏钱的人们而言，买进标普500可能是个好的办法，就当存银行了，利息还高。（2024-01-06）

\BookQuestion{网友：您觉得投资伯克希尔最大的风险在哪里？}

我有不少伯克希尔，不是很懂你的问题。我个人认为伯克希尔大概率可以打败通胀，大于51\%的机会可以打败标普500，以及大概率会打败99\%以上说自己不喜欢巴菲特的人们。

我觉得对多数人而言，买标普500可能是比较好的选择，因为伯克希尔对多数人毕竟还是“一家”公司，有波动的时候可能会更影响心情。虽然我认为伯克希尔的长期表现会好过标普500，但也不会好很多的，而拿着标普500对多数很难理解伯克希尔的人来说心情会好一些。（2022-03-21）

被通胀抢，好过在股市亏。（2010-10-23）

我发现价值投资对很多人来讲就像对面那座大山，看起来很近，实际上可能还很远。（2011-10-21）

假装会游泳是很危险的，假装会投资也一样。不过很多人不很理解假装会投资的危险到底在哪里，直到他们理解了危险。（2023-01-10）

炒股最危险的日子是礼拜一，其次分别是礼拜五，礼拜四，礼拜二，礼拜三。类似的话好像是海明威说的？（2019-05-24）

我理解的所谓投资其实和“市场”（短期）没关系，所以没有谁比谁聪明的问题，但有每个人了解的东西不一样的概念。我认为的投机就是认为自己可以比市场快一步。我认为的投资和市场没关系（意思就是别人短期怎么看没关系）。如果一个投机分子认为市场比自己聪明的话，从一开始他就不会介入（来小玩玩的不算哈），当然最后大部分人确实证明了市场更聪明。从长期的角度看，市场是绝顶聪明的！其实很多“购买并持有”的人们实际上就是投机分子，因为他们自己其实不知道自己买的是什么，所以买了以后需要整天问别人怎么看，在被套时就拿出“长期投资”来安慰自己。要知道，没什么比买错股票并长期持有伤害更大的了。（2011-10-29）

\BookQuestion{网友：那些主流理念究竟给我们带来了什么？}

只想说一点：“长期投资”并不自动等于价值投资，但价值投资一般都是长期投资。（2012-03-14）

价值投资并不是盲目地买进并长期持有。如果是好公司才值得在好价格时买进并长期持有。当然，好价格并不一定是抄底的价格（如果是当然好），好价格指的是10年后回头看：哇塞，那真是好价格。（2010-04-30）

也许“长期投资”这个说法有点问题，叫“投资长期”似乎不容易误解。（2024-08-10）

\BookQuestion{网友：“理解复利的魔力和获得它的困难是理解许多事情的核心和灵魂。”芒格这句话怎么理解？}

不做不对的事情，从而少犯错。对于具有把事情做对的能力的人而言，时间长了就明白什么是复利了。（2013-06-04）

\BookQuestion{网友：很好奇您作为一个优秀的投资人，有没有犯过错？}

我投资也犯过错，以后也还会犯。以前犯错可能是因为做了错的事情，也可能是在做对的事情过程中犯的错。以后我犯的错大概绝大多数会是把事情做对的能力上的问题。（2019-03-14）

\BookQuestion{网友：如果一家公司为上市融资获得高估值而财务造假，导致董事长被捕入狱。但这家公司所处行业毛利率在30\%以上，而且刚签了大合同。在公司出现大跌时，我买入并耐心长期持有，打电话询问公司主营业务的进展，新董事长会在公司管理上做出调整来防范再次出现危机。我这样做算不算一个投机者？}

如果你真打算持有10年的话就不算，希望你能睡好觉。（2011-03-22）

\BookQuestion{网友：这次公司出了事，曝光得越多越好，矫枉过正希望能置之死地而后生。}

我发现做错的事情的人和公司往往在侥幸活过来后还会继续做错的事情，就有点像骗子被抓以后往往会认为是自己骗术不够好而总结如何提高骗术一样。（2011-03-24）

\BookQuestion{网友：价值投资“术”的层面怎么看？}

我个人以为价值投资的术就是如何研究企业未来现金流（折现）的问题。巴菲特的护城河和好的管理团队等都属于术的范畴。其实这个术是很难从书本上学到的，学校里一般也没人可以教，大概只能自己慢慢悟了。（2011-01-05）

对了解投资而言，读书多确实未必有帮助，因为关于投资的好书很少，有投资鉴别力的读者可能也不是很多。（2022-10-26）

\BookQuestion{网友：最想学的就是您说的“把事情做对”，并坚持下去。}

把事情做对是要靠学习的，在实践中学习，在犯错误中学习。不做永远都不犯错，但永远学不会。（2010-03-27）

把事情做对确实有高低之分。每个人都有自己的学习曲线，每个人都不可能不犯错误。据说在任何方面花够一万个小时后，大多数人都能成为专家。（2015-03-25）

% 来源：EPUB part043.xhtml
\subsection{快即是慢}

\BookQuestion{问：您竞拍巴菲特午餐时用的ID是“Fast is slow”，您能解释一下是什么意思吗？}

太快了就等于慢，其实就相当于我们中国人说的“欲速不达”。这是我多年以来做企业、做投资的一个最基本的哲学。做很多事情你必须要认真去做，要有耐心，要有坚持，不要老想着“多快好省”，走捷径的办法往往都是走不远的，这是一个很基本的逻辑。（2006-07-07）

我一直说“Fast is slow”，原意是“欲速不达”，不知道为什么“后人”会说成“慢就是快”。欲速不达反过来并不是慢就是快，逻辑上并不成立。（2025-03-02）

我觉得设定投资目标是件蛮愚蠢的事情。我的意思是：当你有了每年增长多少百分比的目标之后，你早晚会做出很愚蠢的事情。（2023-03-01）

一旦心里有了要求回报率，动作就可能变形，因为人们是有可能为了达到这个要求回报率去铤而走险的。还是那句老话：对投资而言过程正确最重要，心里有个要求回报率大概率是会降低你的回报率的。也许可以有个“可接受回报率”，就是低于这个可接受回报率的投资就不碰了。我觉得这个可接受回报率的底线其实就是国债利率。（2024-07-18）

\BookQuestion{网友：你在网易、苹果等公司的投资收益非常高，但关于个股选择的机会成本，却只跟美国国债（收益很低）比较，而不是跟标普等指数比较。为什么？}

其实跟什么比并不能改变你投资的原则。我喜欢选个标准比较低一点的，满足感比较好。要长期战胜标普非常难，虽然我投资以来的实际成绩总的来说是好过标普的，但我宁愿将其归之于运气，因为我无法知道未来我还能不能做到。（2022-01-30）

我个人认为，投资最重要的动机是保值，保持购买力不下降。其他都是附带的。多数人可能没有意识到保值不是件容易的事情。看下M2的数字也许能明白点？（2020-10-26）

投资是抵抗通胀的最好办法，不是最好办法之一。所以没有闲钱投资是件痛苦的事情。（2012-06-26）

做固定收益是一定会输给通胀的，所以长期而言是错的。但是，如果你实在是找不到愿意投的公司，那投固定收益也是一个选择，肯定比拿着现金放在保险柜里好。（2022-01-22）

\BookQuestion{网友：在10年投资中您都犯过哪些错误？补了哪些弱项？感到最难的是什么？}

这些年犯的错误就那几个：做了自己其实不懂的东西，借过钱，做过空。所有的错误都和这些有关。以后不会再做空和借钱，但可能还会有机会做以为自己懂了其实还是没搞懂的东西。所以犯错率大幅度下降了。最难的就是什么都不做。当手里有很多现金的时候就很容易理解为什么这么说。巴菲特说他最容易犯错的时候就是手里有很多现金的时候，我也是。（2010-03-08）

\BookQuestion{网友：当手头有闲置的现金又找不到合适的股票买入，内心里觉得钱如果没能“花”出去就是在贬值，我会着急。换作是您，您会怎么做呢？}

我会着急，但是，着急时尽量不要决策。决策要建立在理性的基础上。（2019-04-07）

其实茅台2500元左右时我也想过卖掉的问题，最后还是放弃了。因为卖掉后真不知道拿着现金干什么，利息还那么低，人会焦虑的，不合算。

清仓茅台不是问题，问题是拿着那笔钱都干嘛了。绝大部分人很可能是躲过了和尚但没躲过庙？（2024-01-03）

其实你卖了茅台去换腾讯是个过得去的决定，因为最后你也不会亏到哪里去。很多人卖了茅台然后买了比如“假布斯”之类再也回不来的东西才是真的惨了。（2024-05-26）

\BookQuestion{网友：高估时卖掉，若有定力存在银行一直等到下一次低估当然最好，但有几个投资者有如此定力？}

需要“定力”的东西一般来说不会是啥让人开心的东西，甚至可能是让人睡不好觉的东西。

\BookQuestion{网友：大道眼力够强，所以不需要定力。普通人没有大道的眼力，又想理财，还是要选择一条适合自己的路。每个人的悟性能力不同，强求与大道一样的投资方法或许会南辕北辙。}

嗯嗯，尽管找借口吧！对你没啥好处的。

\BookQuestion{网友：本分！老老实实……就不大需要定力。}

关键还是要理性。（2024-05-26）

\BookQuestion{网友：没钱时遇到好公司怎么办？}

钱多多买点，钱少少买点呗。（2020-10-11）
